\subsection*{平面向量}
\begin{enumerate}
    \item 已知 $\bigtriangleup {OAB}$ 中, $\vv{OA} = \vv{a},\vv{OB} = \vv{b}$ ,满足 $\left| \vv{a}\right|  = \left| \vv{b}\right|  = \left| {\vv{a} - \vv{b}}\right|  = 3$ ,求 $\left| {\vv{a} + \vv{b}}\right|$ 与 $\bigtriangleup {OAB}$ 面积.
    \begin{jieda}
        由已知得 $\left| \vv{OA}\right|  = \left| \vv{OB}\right|$ ,以 ${OA},{OB}$ 为邻边作平行四边形 ${OACB}$ ,则可知其为菱形, 且 $\vv{OC} = \vv{a} + \vv{b},\vv{BA} = \vv{a} - \vv{b}$ ,由于 $\left| \vv{a}\right|  = \left| \vv{b}\right|  = \left| {\vv{a} - \vv{b}}\right|  = 3$ ,则 $\left| {OA}\right|  = \left| {OB}\right|  = \left| {BA}\right|  = 3$，根据几何关系可知$\left| {\vv{a} + \vv{b}}\right|  = \left| {OC}\right|  = 3\sqrt{3}$，${S}_{\bigtriangleup {OAB}} = \dfrac{1}{2}\left| {OA}\right|  \cdot  \left| {OB}\right| \sin \dfrac{\pi }{3} = \dfrac{\sqrt{3}}{4} \times  {3}^{2} = \dfrac{9\sqrt{3}}{4}$
        
    
        \begin{center}
        \includegraphics[width=0.3\textwidth]{bodyxb/src/bo_d5dqkaref24c73bjk040_0_240_1570_372_250_0.jpg}
        \end{center}

    \end{jieda}
    \item 设平面上有四个互异的点$A,B,C,D$，且$(\vv{DB} + \vv{DC} -2\vv{DA}) \cdot (\vv{AB} - \vv{AC}) = 0$，则$\triangle ABC$的形状是\dotfill{(\kh{B})}
    \xx{直角三角形}{等腰三角形}{等腰直角三角形}{等边三角形}
    \begin{jieda}
        $(\vv{DB} + \vv{DC} -2\vv{DA}) \cdot (\vv{AB} - \vv{AC}) = 0$即$(\vv{AB} + \vv{AC}) \cdot \vv{CB} = 0$，即$BC$边中线与$BC$垂直，根据几何关系可知$AB = AC$，即$\triangle ABC$为等腰三角形.
    \end{jieda}
    \item 已知$A,O,B$是不共线的三点，且$C,D$满足$\vv{OC} = 2 \vv{OA}, \vv{OD} = 3 \vv{OB}$，$AD$与$BC$交于点$M$，设$\vv{OM} = x \vv{OA} + y \vv{OB}$，试求$x, y$.
    \begin{jieda}
        $\vv{OM} = x \vv{OA} + y \vv{OB} \Rightarrow \left \{ \begin{array}{l}
                \vv{OM} = \dfrac{x}{2}\vv{OC} + y\vv{OB}  \\
                \vv{OM} = x \vv{OA} + \dfrac{y}{3}\vv{OD} 
        \end{array} \right.$，因为$A, M, D$共线，$C, M, B$共线，根据等和线结论有$\left \{ \begin{array}{l}
                \dfrac{x}{2} + y = 1  \\
                x + \dfrac{y}{3} = 1 
        \end{array} \right.$，解得$x = \dfrac{4}{5}, y = \dfrac{3}{5}$.
    \end{jieda}

    \item 若向量$\vv{a} = \left(\dfrac{3}{2}, sin\alpha\right), \vv{b} = \left(cos\alpha, \dfrac{1}{3}\right)$，且$\vv{a} // \vv{b}$，则锐角$\alpha = $\blkda{$\dfrac{\pi}{4}$}
    \begin{jieda}
        根据平面向量平行的充要条件的坐标表示有$sin\alpha \cdot cos\alpha = \dfrac{1}{2}$，故$sin2\alpha = 1$，因此$\alpha = \dfrac{\pi}{4}$
    \end{jieda}

    \item 设$\vv{a}, \vv{b}, \vv{c}$是任意的非零向量，且两两不共线，则下列命题中正确的是\blkda{\ding{173} \ding{175}}:
    \begin{dingautolist}{172}
        \item $(\vv{a} \cdot \vv{b}) \cdot \vv{c} - (\vv{c} \cdot \vv{a}) \cdot \vv{b} = \vv{0}$
        \item $|\vv{a}| - |\vv{b}| < |\vv{a} - \vv{b}|$
        \item $(\vv{b} \cdot \vv{c}) \cdot \vv{a} - (\vv{a} \cdot \vv{c}) \cdot \vv{b}$不与$\vv{c}$垂直
        \item $(3\vv{a} + 2\vv{b}) \cdot (3\vv{a} - 2\vv{b}) = 9|\vv{a}|^2 - 4|\vv{b}|^2$
    \end{dingautolist}


    \item 已知向量 $\vv{AB} = \left( {k,1}\right) ,\vv{AC} = \left( {1,0}\right)$ ， $\bigtriangleup  {ABC}$ 是直角三角形，则 $k =$ \blkda{0 或 1}
    \begin{jieda}
        因为 $\vv{AB} = \left( {k,1}\right) ,\vv{AC} = \left( {1,0}\right)$ ,所以 $\vv{CB} = \vv{AB} - \vv{AC} = \left( {k - 1,1}\right)$ ,
        
        若 $\vv{AB} \bot  \vv{AC}$ ,则 $\vv{AB} \cdot  \vv{AC} = k = 0$ ,所以 $k = 0$ ,符合题意;
        
        若 $\vv{AB} \bot  \vv{CB}$ ,则 $\vv{AB} \cdot  \vv{CB} = k\left( {k - 1}\right)  + 1 = 0$ ,此方程无解;
        
        若 $\vv{AC} \bot  \vv{CB}$ ,则 $\vv{AC} \cdot  \vv{CB} = k - 1 = 0$ ,所以 $k = 1$ ,符合题意;
        
        综上, $k = 0$ 或 1 .
    \end{jieda}

    \item 如图所示,已知 $\bigtriangleup  {ABC}$ 中，点 $P,Q,R$ 依次是边 ${BC}$ 上的三个四等分点，若 ${BC} = 8$ ， $\vv{AP} \cdot  \vv{AR} = {20}$ ，则 $\vv{AB} \cdot  \vv{AC} =$ \blkda{8}.
    
    \begin{center}
    \includegraphics[width=0.2\textwidth]{bodyxb/src/3_241_1345_346_180_0.jpg}
    \end{center}
    \begin{jieda}
        $20 = \vv{AP} \cdot  \vv{AR} = |\vv{AQ}|^2 - |\vv{PQ}|^2 = |\vv{AQ}|^2 - 4$

        $\vv{AB} \cdot  \vv{AC} = |\vv{AQ}|^2 - |\vv{BQ}|^2 = |\vv{AQ}|^2 - 16 = 8$
    \end{jieda}

    \item 已知$\vv{e_1}, \vv{e_2}$是两个夹角为$30^{\circ}$的单位向量，且$\vv{a} \cdot \vv{e_1} + 2 = 0$，$|\vv{b} - \vv{e_2}| = \dfrac{1}{2}$，则$|3\vv{a} - 2\vv{b}|$的最小值为\blkda{$5+\sqrt{3}$}
    \begin{jieda}
        首先以$\vv{e_1}$为正方向，$\vv{e_2}$在第一象限，建立平面直角坐标系，由此$\vv{e_1} = (1, 0), \vv{e_2} = (\dfrac{\sqrt{3}}{2}, \dfrac{1}{2})$，设$\vv{a} = (m, n), \vv{b} = (p, q)$，则题干条件可以转化为：$m = -2$，$(p-\dfrac{\sqrt{3}}{2})^2+(q-\dfrac{1}{2})^2 = \dfrac{1}{4}$ \\
        设原点为$O$，故$\vv{OA} = \vv{a}$的终点$A$在直线$x = -2$上，$\vv{OB} = \vv{b}$的终点$B$在以$(\dfrac{\sqrt{3}}{2}, \dfrac{1}{2})$为圆心，$\dfrac{1}{2}$为半径的圆上，进一步转化有$\vv{OC} = 3\vv{a}$的终点$C$在直线$x = -6$上，$\vv{OD} = 2\vv{b}$的终点$D$在以$(\sqrt{3}, 1)$为圆心，1为半径的圆上。\\
        因此$|3\vv{a} - 2\vv{b}|$即为以圆$(x-\sqrt{3})^2+(y-1)^2 = 1$上一点为起点，以直线$x = -6$上一点为终点的向量的模，故其最小值为$5+\sqrt{3}$
    \end{jieda} 

    \item 已知平面内的向量$\vv{a}, \vv{b}, \vv{c}, \vv{e}$满足$|\vv{a}| = 4, |\vv{e}| = 1, \langle\vv{a}, \vv{e}\rangle = \dfrac{2\pi}{3}, |\vv{b} - \vv{a}| = 1$，且有$|\vv{c} - \lambda \vv{e}| \geq |\vv{c} - 2\vv{e}|$对任意$\lambda \in \mathbf{R}$恒成立，则$|\vv{c} - \vv{b}|$的最小值为\blkda{3}.
    
    \item 已知圆$O$的半径为1，$PA, PB$为该圆的两条切线，$A, B$为两切点，那么$\vv{PA} \cdot \vv{PB}$的最小值是\blkda[2]{$-3 + 2\sqrt{2}$}
    \begin{jieda}
        $\vv{PA} = \vv{PO} + \vv{OA}, \vv{PB} = \vv{PO} + \vv{OB}, \vv{PA} \cdot \vv{PB} = |\vv{PO}|^2 + \vv{PO} \cdot (\vv{OA} + \vv{OB}) + \vv{OA} \cdot \vv{OB}$ \\
        若设$|\vv{OP}| = x, \angle POA = \alpha$，易知$cos\alpha = \dfrac{1}{x}$，则$\vv{PA} \cdot \vv{PB} = x^2 - 2x\cdot cos\alpha + cos2\alpha = x^2 - 2 + (\dfrac{2}{x^2} - 1) = -3 + x^2 + \dfrac{2}{x^2} \geq -3 + 2\sqrt{2}$，当$x = \sqrt[4]{2}$时取等.
    \end{jieda}

    \item 在 $\bigtriangleup {ABC}$ 中,角 $A,B,C$ 所对的边分别是 $a,b,c$ ,若 ${\sin }^{2}B + {\sin }^{2}C = {\sin }^{2}A + \sin B\sin C$ ，且 $\vv{AC} \cdot  \vv{AB} = 4$ ， $\bigtriangleup  {ABC}$ 的面积 $S$ 为\blkda{$2\sqrt{3}$}.
    \begin{jieda}
        由正弦定理知， ${b}^{2} + {c}^{2} = {a}^{2} + {bc}$ ，
        
        由余弦定理知, ${b}^{2} + {c}^{2} - {a}^{2} = {2bc}\cos A = {bc}$, 则 $\cos A = \dfrac{1}{2},A = \dfrac{\pi }{3}$
        
        又 $\vv{AC} \cdot  \vv{AB} = \left| \vv{AC}\right|  \cdot  \left| \vv{AB}\right| \cos A = \dfrac{1}{2}{bc} = 4$
        
        则三角形面积 $S = \dfrac{1}{2}{bc}\sin A = 4 \times  \dfrac{\sqrt{3}}{2} = 2\sqrt{3}$
    \end{jieda}

    \item 已知$A,B,C$是平面上不共线上三点, 动点$P$满足$\vv{OP}=\dfrac{1}{3}((1-\lambda)\vv{OA}+(1-\lambda)\vv{OB}+(1+2\lambda)\vv{OC})(\lambda\zr,\lambda\neq 0)$，则$P$的轨迹一定通过$\triangle{ABC}$的\dotfill(\kh{D})
    \xx{内心}{垂心}{重心}{$AB$边的中点}
    \begin{jieda}
        $\vv{OP}=\dfrac{1}{3}((1-\lambda)\vv{OA}+(1-\lambda)\vv{OB}+(1+2\lambda)\vv{OC}) = \dfrac{1}{3}\vv{OA} + \dfrac{1}{3}\vv{OB} + \dfrac{1}{3} \vv{OC} + \dfrac{\lambda}{3}(\vv{AC} + \vv{BC})$\\
        设$G$为重心，$\vv{OP} = \vv{OG} + \dfrac{\lambda}{3}(\vv{AC} + \vv{BC})$，即$\vv{GP} = -\lambda \vv{CG}$,由于$\lambda \neq 0$，所以选D.
    \end{jieda}

    \item 已知 $\bigtriangleup  {ABC}$ 三边的垂直平分线交于点 $O$ ,且 $A{B}^{2} + A{C}^{2} - {2AC} = 0$ ,则 $\vv{BC} \cdot  \vv{AO}$ 的取值范围是\blkda{$\left[ -\dfrac{1}{4},2\right)$}
    \begin{jieda}
        $A{B}^{2} + A{C}^{2} - {2AC} = 0$即$c^2+b^2-2b = 0$

        $\vv{BC} \cdot  \vv{AO} = \vv{AO} \cdot (\vv{AC} - \vv{AB}) = \dfrac{b^2 - c^2}{2} = b^2 - b$

        $2b - b^2 > 0$，故$b \in (0, 2)$，因此$b^2 - b \in \left[ -\dfrac{1}{4},2\right)$
    \end{jieda}

    \item 在 $\bigtriangleup  {ABC}$ 中，角 $A,B,C$ 的对边分别为 $a,b,c$ ，已知 $\dfrac{a}{b} = \dfrac{\sqrt{3}}{3}\sin C + \cos C$ .
    \begin{enumerate}
        \item 求角 $B$
        \item 若 $D$ 是 $\bigtriangleup  {ABC}$ 边 ${AC}$ 上的一点，且满足 $\dfrac{\vv{BA} \cdot  \vv{BD}}{\left| \vv{BA}\right| } = \dfrac{\vv{BD} \cdot  \vv{BC}}{\left| \vv{BC}\right| }$ ， ${9a} + {4c} = {25}$ ，求 ${BD}$ 的最大值.
    \end{enumerate}
    \begin{jieda}
        \begin{enumerate}
            \item 因为 $\dfrac{a}{b} = \dfrac{\sqrt{3}}{3}\sin C + \cos C$ ,即 $a = \dfrac{\sqrt{3}}{3}b\sin C + b\cos C$ ,
            
            由正弦定理可得 $\sin A = \dfrac{\sqrt{3}}{3}\sin B\sin C + \sin B\cos C$ ,
            
            且 $\sin A = \sin \left( {B + C}\right)  = \sin B\cos C + \cos B\sin C$ ,
            
            即 $\sin B\cos C + \cos B\sin C = \dfrac{\sqrt{3}}{3}\sin B\sin C + \sin B\cos C$ ,可得 $\cos B\sin C = \dfrac{\sqrt{3}}{3}\sin B\sin C$ ,
            
            且 $C \in  \left( {0,\pi }\right)$ ,则 $\sin C \neq  0$ ,可得 $\tan B = \sqrt{3}$ ,
            
            又因为 $0 < B < \pi$ ,所以 $B = \dfrac{\pi }{3}$ .

            \item 因为 $\dfrac{\vv{BA} \cdot  \vv{BD}}{\left| \vv{BA}\right| } = \dfrac{\vv{BD} \cdot  \vv{BC}}{\left| \vv{BC}\right| }$ ，即 $\dfrac{\vv{BA} \cdot  \vv{BD}}{\left| \vv{BA}\right| \left| \vv{BD}\right| } = \dfrac{\vv{BD} \cdot  \vv{BC}}{\left| \vv{BC}\right| \left| \vv{BD}\right| }$ ，
            
            可得 $\cos \angle {ABD} = \cos \angle {CBD}$ ，即 $\angle {ABD} = \angle {CBD}$ ，
            
            可知 ${BD}$ 平分 $\angle {ABC}$ ，则 $\angle {ABD} = {CBD} = \dfrac{\pi }{6}$ ，
            
            因为 ${S}_{\bigtriangleup {ABC}} = {S}_{\bigtriangleup {ABD}} + {S}_{\bigtriangleup {BCD}}$ ，
            
            即 $\dfrac{1}{2}{ac} \times  \dfrac{\sqrt{3}}{2} = \dfrac{1}{2}{BD} \times  a \times  \dfrac{1}{2} + \dfrac{1}{2}{BD} \times  c \times  \dfrac{1}{2}$ ，整理可得 $\dfrac{\sqrt{3}}{BD} = \dfrac{1}{a} + \dfrac{1}{c}$ ，
            
            又因为 ${9a} + {4c} = {25}$ ,
            
            则 $\dfrac{\sqrt{3}}{BD} = \dfrac{1}{25}\left( {{9a} + {4c}}\right) \left( {\dfrac{1}{a} + \dfrac{1}{c}}\right)  = \dfrac{1}{25}\left( {{13} + \dfrac{4c}{a} + \dfrac{9a}{c}}\right)  \geq  \dfrac{1}{25}\left( {{13} + 2\sqrt{\dfrac{4c}{a} \cdot  \dfrac{9a}{c}}}\right)  = 1$ ,
            
            当且仅当 $\dfrac{4c}{a} = \dfrac{9a}{c}$ ,即 $a = \dfrac{5}{3},c = \dfrac{5}{2}$ 时取等号,
            
            可得 ${BD} \leq  \sqrt{3}$ ，所以 ${BD}$ 的最大值为 $\sqrt{3}$ .
        \end{enumerate}
    \end{jieda}
    \item 已知 $a,b,c$ 分别是 $\bigtriangleup  {ABC}$ 三个内角 $A,B,C$ 的对边，且 ${b}^{2} = {\left( a - c\right) }^{2} + 8$ ， ${2b}\sin \left( {A + \dfrac{\pi }{6}}\right)  = a + c$ ,若点 $P$ 为 $\bigtriangleup  {ABC}$ 的费马点，则 $\vv{PA} \cdot  \vv{PB} + \vv{PB} \cdot  \vv{PC} + \vv{PA} \cdot  \vv{PC} =$\blkda{-4}
    
    (注: 已知一个三角形, 求作一点, 使其与这个三角形的三个顶点的距离之和最小的答案是: 当三角形的三个角均小于 ${120}^{ \circ  }$ 时,所求的点为三角形的正等角中心,即该点与三角形的三个顶点的连线两两成角 ${120}^{ \circ  }$ : 当三角形有一内角大于或等于 ${120}^{ \circ  }$ 时,所求点为三角形最大内角的顶点.所求的点称为费马点)
    \begin{jieda}
        已知 ${2b}\sin \left( {A + \dfrac{\pi }{6}}\right)  = a + c$ ,由正弦定理得 $2\sin B\sin \left( {A + \dfrac{\pi }{6}}\right)  = \sin A + \sin C$ ,
        
        由 $\sin \left( {A + \dfrac{\pi }{6}}\right)  = \dfrac{\sqrt{3}}{2}\sin A + \dfrac{1}{2}\cos A,\sin C = \sin \left( {A + B}\right)  = \sin A\cos B + \cos A\sin B$ ,
        
        则有 $\sqrt{3}\sin B\sin A + \sin B\cos A = \sin A + \sin A\cos B + \cos A\sin B$ ,
        
        即 $\sqrt{3}\sin B\sin A = \sin A + \sin A\cos B$ ,
        
        $A \in  \left( {0,\pi }\right) ,\sin A \neq  0$ ,有 $\sqrt{3}\sin B = 1 + \cos B$ ,得 $\sin \left( {B - \dfrac{\pi }{6}}\right)  = \dfrac{1}{2}$ ,
        
        因为 $B \in  \left( {0,\pi }\right)$ ,所以 $B - \dfrac{\pi }{6} \in  \left( {\dfrac{\pi }{6},\dfrac{5\pi }{6}}\right)$ ,所以 $B - \dfrac{\pi }{6} = \dfrac{\pi }{6}$ ,所以 $B = \dfrac{\pi }{3}$ .
        
        由三角形内角和性质知: $\bigtriangleup {ABC}$ 内角均小于 ${120}^{ \circ  }$ ,
        
        结合题设易知: $P$ 点一定在三角形的内部,
        
        再由余弦定理知, $\cos B = \dfrac{{a}^{2} + {c}^{2} - {b}^{2}}{2ac} = \dfrac{1}{2}$ ,
        
        又因为 ${b}^{2} = {\left( a - c\right) }^{2} + 8$ ,所以 ${ac} = 8$ ,
        
        所以 ${S}_{\bigtriangleup {ABC}} = \dfrac{1}{2}\left| {PA}\right|  \cdot  \left| {PB}\right| \sin \dfrac{2\pi }{3} + \dfrac{1}{2}\left| {PB}\right|  \cdot  \left| {PC}\right| \sin \dfrac{2\pi }{3} + \dfrac{1}{2}\left| {PA}\right|  \cdot  \left| {PC}\right| \sin \dfrac{2\pi }{3}$
        
        $= \dfrac{1}{2}{ac}\sin B = \dfrac{1}{2} \times  8 \times  \sin \dfrac{\pi }{3},$
        
        所以 $\left| {PA}\right|  \cdot  \left| {PB}\right|  + \left| {PB}\right|  \cdot  \left| {PC}\right|  + \left| {PA}\right|  \cdot  \left| {PC}\right|  = 8$ .
        
        由 $\left| {PA}\right|  \cdot  \left| {PB}\right|  + \left| {PB}\right|  \cdot  \left| {PC}\right|  + \left| {PA}\right|  \cdot  \left| {PC}\right|  = 8$ ,等号左右两边同时乘以 $\cos \dfrac{2\pi }{3}$ 可得:
        
        $\left| {PA}\right|  \cdot  \left| {PB}\right| \cos \dfrac{2\pi }{3} + \left| {PB}\right|  \cdot  \left| {PC}\right| \cos \dfrac{2\pi }{3} + \left| {PA}\right|  \cdot  \left| {PC}\right| \cos \dfrac{2\pi }{3} = 8 \times  \cos \dfrac{2\pi }{3}$ ,
        
        则 $\vv{PA} \cdot  \vv{PB} + \vv{PB} \cdot  \vv{PC} + \vv{PA} \cdot  \vv{PC} = 8 \times  \cos \dfrac{2\pi }{3} =  - 4$ .
    \end{jieda}
\end{enumerate}